\documentclass[11pt]{article}
\usepackage[margin=1in]{geometry}
\usepackage{amsmath,amssymb,amsthm,mathtools}
\usepackage{graphicx}
\usepackage{microtype}
\usepackage[hidelinks]{hyperref}

\newtheorem{theorem}{Theorem}
\newtheorem{lemma}{Lemma}
\newcommand{\R}{\mathbb{R}}
\newcommand{\N}{\mathbb{N}}
\newcommand{\e}{\mathrm{e}}

\title{A finite-codimensional obstruction for the mixed-weight\\
heat-decay question of Bonet--Lusky--Taskinen}
\author{Automated mathematics run \texttt{fa\_banach\_001}\\
Lane 11, model GPT5.6}
\date{9 August 2026}

\begin{document}
\maketitle

\begin{center}
\fbox{\parbox{0.9\linewidth}{
\textbf{Status.} Candidate counterexample / complete negative answer, likely
valid, pending human review.
}}
\end{center}

\section{Source question}

Bonet, Lusky, and Taskinen study Schauder bases whose finite-codimensional
tails force arbitrarily fast decay for the heat equation \cite{BLT}.  Their
Question $3^\circ$ asks whether their higher-dimensional conclusion can hold
when the weight is fast in only one coordinate:
\[
 w(x)=\e^{|x_1|}\prod_{j=2}^{N}(1+|x_j|)^2,
 \qquad N>1.
\]
The question appears on page 14 of the source PDF, as equation (4.19).

\begin{figure}[ht]
  \centering
  \includegraphics[width=\textwidth]{figures/open_problem_crop.png}
  \caption{The source discussion and Question $3^\circ$ on page 14 of
  arXiv:1802.01902.}
\end{figure}

More explicitly, the requested conclusion would give a Schauder basis
$(e_n)$ of $L^1_w(\R^N)$ such that, for every $m\in\N$, some closed basis tail
$E_m$ satisfies
\begin{equation}\label{eq:desired}
 \|\e^{t\Delta}f\|_\infty
 \le C_m t^{-m}\|f\|_{1,w}
 \qquad (f\in E_m,\ t\ge 1).
\end{equation}
Every such $E_m$ has finite codimension.  We prove that no
finite-codimensional subspace can have decay faster than
$t^{-N/2-1}$ in operator norm.

\section{Proof intuition}

Let $E$ have codimension $d$.  At time $t$, place $d+1$ copies of one fixed
compactly supported bump along the $x_2$-axis, with spacing $L\sqrt t$.  Their
supports are disjoint, and their span must meet $E$ nontrivially.  The mixed
weight of the $j$th translated bump is only $O(j^2t)$, because the weight is
quadratic in $x_2$.

Normalize a nonzero combination in that intersection.  At the center of the
bump having the largest coefficient after weighted normalization, its own
heat contribution is $\asymp t^{-N/2}$.  All other contributions are smaller
by fixed Gaussian factors $\exp(-cL^2|j-k|^2)$, so choosing $L$ large prevents
cancellation.  Weighted normalization costs one further factor $t^{-1}$.
This produces the unavoidable lower bound $t^{-N/2-1}$.

\section{The obstruction theorem}

Let
\[
 X=L^1_w(\R^N),\qquad
 \|f\|_{1,w}=\int_{\R^N}|f(x)|w(x)\,dx,
\]
with the mixed weight above.  Write the heat kernel in the form
\[
 (H_tf)(x)=(\e^{t\Delta}f)(x)
 =\gamma_N t^{-N/2}\int_{\R^N}
   \exp\!\left(-\frac{|x-y|^2}{4t}\right)f(y)\,dy,
\]
where the positive normalization constant $\gamma_N$ is immaterial.

\begin{theorem}\label{thm:main}
Let $N>1$ and let $E\subset X$ be a closed subspace of finite codimension.
There are constants $c_E>0$ and $T_E\ge1$ such that
\[
 \|H_t|_E\|_{X\to L^\infty(\R^N)}
 \ge c_E t^{-N/2-1}
 \qquad (t\ge T_E).
\]
Consequently, no Schauder basis of $X$ has the arbitrary tail-decay property
\eqref{eq:desired}.  Thus Question $3^\circ$ has a negative answer.
\end{theorem}

\begin{proof}
Fix a nonnegative, nonzero function
$\varphi\in C_c^\infty(\R^N)$ whose support lies in the ball
$B(0,\rho)$.  Let
\[
 d=\dim(X/E),\qquad D=d+1.
\]
Choose a constant $L>0$, to be fixed below, and for $t\ge1$ define centers
and weighted normalizing constants by
\[
 z_j(t)=jL\sqrt t\,e_2,
 \qquad
 a_{j,t}=\int_{\R^N}\varphi(u)w(z_j(t)+u)\,du,
 \qquad 1\le j\le D.
\]
Set
\[
 g_{j,t}(x)=\frac{\varphi(x-z_j(t))}{a_{j,t}}.
\]
For all sufficiently large $t$, the supports of these $D$ functions are
pairwise disjoint, and $\|g_{j,t}\|_{1,w}=1$.

The $D$-dimensional space spanned by the $g_{j,t}$ meets $E$ nontrivially,
because $X/E$ has dimension $d<D$.  Hence there are scalars $c_1,\ldots,c_D$,
not all zero, such that
\[
 f_t=\sum_{j=1}^{D}c_jg_{j,t}\in E.
\]
After rescaling, assume $\sum_j|c_j|=1$.  Disjointness of the supports then
gives the exact identity
\[
 \|f_t\|_{1,w}=\sum_{j=1}^{D}|c_j|=1.
\]

Put $q_j=|c_j|/a_{j,t}$ and choose $k$ with
$q_k=\max_jq_j$.  On the fixed support of $\varphi$, all weight factors other
than the $x_2$ factor are bounded above by constants independent of $j$ and
$t$, while
\[
 (1+|jL\sqrt t+u_2|)^2\le C_Lj^2t
 \qquad (t\ge1,\ |u_2|\le\rho).
\]
It follows that
\begin{equation}\label{eq:aj}
 a_{j,t}\le Cj^2t,
 \qquad
 \sum_{j=1}^{D}a_{j,t}\le C_Et.
\end{equation}
Since $1=\sum_j|c_j|=\sum_jq_ja_{j,t}$, we obtain
\begin{equation}\label{eq:q}
 q_k\ge \frac{1}{\sum_{j=1}^{D}a_{j,t}}
 \ge C_E^{-1}t^{-1}.
\end{equation}

We now evaluate $H_tf_t$ at $z_k(t)$.  The diagonal integral satisfies
\begin{align}
 I_{k,t}
 &:=\gamma_Nt^{-N/2}\int
 \exp\!\left(-\frac{|u|^2}{4t}\right)\varphi(u)\,du \notag\\
 &\ge \gamma_N\e^{-\rho^2/4}\|\varphi\|_1t^{-N/2}
 =:A_0t^{-N/2}.                                      \label{eq:diag}
\end{align}
For $j\ne k$ and all sufficiently large $t$,
\[
 |z_k(t)-z_j(t)-u|
 \ge L|j-k|\sqrt t-\rho
 \ge \frac L2|j-k|\sqrt t
 \qquad (u\in\operatorname{supp}\varphi),
\]
and therefore the corresponding off-diagonal integral is at most
\begin{equation}\label{eq:off}
 I_{j,k,t}
 \le \gamma_N\|\varphi\|_1t^{-N/2}
 \exp\!\left(-\frac{L^2|j-k|^2}{16}\right).
\end{equation}
Choose $L$ once and for all so large that
\[
 2\gamma_N\|\varphi\|_1
 \sum_{r=1}^{\infty}\exp(-L^2r^2/16)\le \frac{A_0}{2}.
\]
Using $q_j\le q_k$, the reverse triangle inequality together with
\eqref{eq:diag}--\eqref{eq:off} yields
\begin{align*}
 |(H_tf_t)(z_k(t))|
 &\ge q_k I_{k,t}-\sum_{j\ne k}q_jI_{j,k,t}\\
 &\ge \frac{A_0}{2}q_kt^{-N/2}\\
 &\ge c_Et^{-N/2-1},
\end{align*}
where the last line uses \eqref{eq:q}.  Since $f_t\in E$ and
$\|f_t\|_{1,w}=1$, this proves the operator-norm lower bound.

Finally, suppose a Schauder basis had the property requested in the source.
Choose an integer $m>N/2+1$ and let $E_m$ be its corresponding closed basis
tail.  This tail has finite codimension, so the lower bound just proved and
\eqref{eq:desired} would imply, for all sufficiently large $t$,
\[
 c_{E_m}t^{-N/2-1}\le C_mt^{-m}.
\]
Multiplication by $t^{N/2+1}$ and passage to $t\to\infty$ gives
$c_{E_m}\le0$, a contradiction.  Hence the requested basis does not exist.
\end{proof}

\section{Verification notes and scope}

The proof establishes more than failure of one tensor-product construction:
it obstructs the desired estimate on every closed finite-codimensional
subspace of $L^1_w(\R^N)$.  The exponent $N/2+1$ comes from ordinary heat
smoothing ($N/2$) plus the quadratic transverse weight (one further power of
$t$ at distance $\sqrt t$).

The following points were checked adversarially.
\begin{enumerate}
\item Translation occurs only in the second coordinate, which exists because
      $N>1$; the exponential $x_1$ factor therefore remains uniformly bounded
      on the translated bump supports.
\item The vector in $E$ may have arbitrary real or complex coefficients.
      Positivity of the Gaussian and the reverse triangle inequality make the
      diagonal-dominance estimate insensitive to their phases.
\item The test vector may depend on $t$, as is appropriate for a lower bound
      on the operator norm of $H_t|_E$.
\item Constants are allowed to depend on $E$ (through $D=d+1$), but never on
      $t$.  This is enough to contradict any faster power bound.
\item No computational assertion or unproved auxiliary lemma enters the
      proof.
\end{enumerate}

\section{Novelty check and limitations}

A bounded search on 9 August 2026 used arXiv, exact-title searches, the author
names, the displayed weight, and the phrases ``fast growing in one
coordinate'' and ``Schauder basis heat equation.''  It located the arXiv
source and author-hosted versions, including a later manuscript that proves a
stronger higher-dimensional theorem for weights fast in every direction.  No
paper or manuscript answering the displayed mixed-weight question was found.
The later theorem does not apply here: the factors $(1+|x_j|)^2$ fail the
all-orders growth assumption in the transverse directions.

This is a bounded, not exhaustive, novelty review.  Human review should focus
on the normalization estimate \eqref{eq:aj}--\eqref{eq:q} and the Gaussian
diagonal-dominance step.  Subject to that review, the argument is a complete
negative answer to the literal Question $3^\circ$.

\begin{thebibliography}{9}
\bibitem{BLT}
J. Bonet, W. Lusky, and J. Taskinen,
\emph{Schauder bases and the decay rate of the heat equation},
arXiv:1802.01902 [math.FA], 2018, Question $3^\circ$, p.~14.
\end{thebibliography}

\end{document}
